\begin{Abstract}
Industrial automations increasingly demands manipulation capability that a single, fixed-workspace
robotic arm cannot provide: components must be perceived reliably even when small and
low-textured, a second arm must be brought in once one arm's reach is exceeded, that second arm
introduces a collision risk no single-arm system faces, and the resulting shared-workspace motion
must be planned, allocated, and, where classical control reaches its limits, learned. This thesis
develops that capability progressively, tracing a single arc from single-arm classical computer
vision to dual-arm reinforcement learning across eight connected studies, each addressing a
research gap the previous study's classical or learned components leave open, and each validated on
a common Kinova Gen3 Lite dual-arm hardware platform.

The first two studies address \textbf{perception} for small, low-textured industrial fasteners.
The first localizes overlapping bolts from RGB-D point clouds using $k$-means segmentation and
Oriented Chamfer matching, achieving bolt sorting from a cluttered bin without misclassification in
Gazebo simulation. The second replaces classical segmentation with a custom-trained DeepLabV3+
model, reaching a mean Intersection over Union of 0.95 against PPLiteSeg and OCRNet baselines, and
fuses the resulting mask with RealSense depth to localize and sort bolts by size on hardware.

The third study addresses \textbf{dual-arm coordination}: two Kinova Gen3 Lite manipulators, each
with its own wrist-mounted eye-in-hand camera rather than a shared external sensor, independently
search a shared workspace, relay a detected target's pose across a calibrated inter-arm transform,
and resolve which arm should act through a runtime, reachability-based role-assignment rule,
validated end-to-end in Gazebo/ROS2 and on hardware with sub-5\,mm, sub-$1.3^\circ$ manipulation
accuracy regardless of which arm is assigned.

The fourth, fifth, and sixth studies address \textbf{safety and motion planning} in the resulting
shared workspace. A kinematics-based framework models manipulator links as line segments,
classifies their configuration as parallel, intersecting, or skew, and computes a continuous,
sensor-free inter-arm distance signal across sixteen collision-critical link pairs in real time.
An improved artificial potential field (iAPF) framework then consumes that signal, coupling a
three-way equilibrium of exponential attraction, continuous inter-arm repulsion, and home-seeking
forces, stabilized by velocity-dependent damping, with a priority-lock coordination mechanism,
and is validated first under a fixed nut-left/bolt-right
object assignment before a distance-based competitive task allocation mechanism replaces that fixed
assignment with a runtime, temporal-window decision, giving an approximately 47\% velocity
improvement over fixed allocation in homogeneous scenarios with zero collisions across all trials.

The seventh and eighth studies replace the classical control layer with \textbf{learned}
controllers exactly where it reaches its limit. A coupled-reward Proximal Policy Optimization (PPO)
controller maps instantaneous Cartesian pose error directly to joint velocity commands for the
Kinova Gen3 Lite's non-spherical wrist, whose offset joint axes break the kinematic decoupling
closed-form inverse kinematics depends on; trained purely on 1.5 million static reachable poses with
no trajectory supervision, it generalizes zero-shot to circular, square, lemniscate, and 3-D
Lissajous trajectories on physical hardware, and is the only method among five classical and
learned baselines that does not degrade from simulation to hardware. That trained policy is then
reused, frozen, as the single-arm prior inside \textbf{ALPHA}, an adaptive hierarchical blending
architecture in which a Multi-Agent PPO coordination policy, trained under Centralized Training
with Decentralized Execution, is combined with the frozen prior through a learned, incrementally
updated gating scalar co-trained with it from timestep zero, achieving 96.5\% reach and 95.0\%
alignment success across 4096 parallel simulation environments and 93.3\% reach with 90.0\%
alignment success over 30 physical hardware trials of cooperative peg-in-hole pre-assembly
alignment, with policy inference at 200\,Hz on a CPU-only platform against the hardware's 20\,Hz
control loop.

Taken together, these eight studies show that the classical and learned layers of dual-arm
industrial manipulation are not alternatives but a single progression, in which each stage supplies
the precondition the next requires: perception must localize a component before either arm can act
on it, a second arm cannot be added safely before it is decided which arm acts, concurrent motion
cannot be trusted without a continuous safety signal, that signal is only useful once a planner
consumes it, and only with a working classical planner in place does it become possible to identify
precisely where a hand-designed control law runs out of headroom and a learned one is warranted.
\\
\textbf{Keywords: -} Dual-arm manipulation, industrial automation, RGB-D perception, template
matching, semantic segmentation, eye-in-hand vision, reachability-based role assignment,
inter-arm collision avoidance, artificial potential fields, competitive task
allocation, differential inverse kinematics, non-spherical wrist manipulators, proximal policy
optimization, multi-agent reinforcement learning, sim-to-real transfer, cooperative pre-assembly
alignment.
\end{Abstract}